% GENERATED FILE. Edit canonical lessons, then run npm run generate:books.
% canonical-source-hash: fnv1a64:23d96d38102a8bbc
% canonical-lessons: MR-C488-rugna, MR-C488-paricarika, MR-C488-vaidya, MR-C488-nidan, MR-C488-upcar

\chapter{At the Clinic}
\label{ch:mr-at-the-clinic}
\hlchaptermodality{488}

\begin{chapteropening}
\textbf{By the end of this chapter:} \emph{I can talk about a patient, a nurse, a physician, a diagnosis and treatment.}

Complete the last lesson of chapter 488: I can talk about a patient, a nurse, a physician, a diagnosis and treatment.
\end{chapteropening}

\section[\texorpdfstring{\mr{रुग्ण} (rugṇa)}{रुग्ण (rugṇa)}]{\mr{रुग्ण} (rugṇa) — a sick person; a patient}

\label{lesson:MR-C488-rugna}



\begin{quote}
Before the new one: say the Marathi for to roast, then the Marathi for to scratch.
\end{quote}

\subsection*{You'll want to know: \mr{रुग्ण}}
\textbf{\mr{रुग्ण}} — \emph{rugṇa} — \textquotedblleft{}a sick person; a patient\textquotedblright{}.

\begin{quote}
\mr{दवाखान्यात} \mr{सकाळपासून} \mr{रुग्णांची} \mr{रांग} \mr{होती}.
\end{quote}

\begin{quote}
\emph{davākhānyāt sakāḷpāsūn rugṇāñcī rāṅg hotī.}
\end{quote}

\begin{quote}
There was a queue of patients at the clinic from the morning.
\end{quote}

Masculine.

As in \mr{रुग्णालय}, a hospital, and \mr{रुग्णवाहिका}, an ambulance.

\begin{cousinweb}
\textbf{\mr{रुग्ण}} (\emph{rugṇa}) is a tatsama, taken from Sanskrit \emph{rugṇa}, \textquotedblleft{}broken, diseased\textquotedblright{} (CDIAL 10760).
\end{cousinweb}

\begin{grammarlens}[title={What you've built}]
One more word for describing what you lived and felt.
\end{grammarlens}

\subsection*{Your turn}
\begin{itemize}
\raggedright
  \item \emph{Say it:} \emph{rugṇa}
  \item \emph{Say it:} \emph{rugṇa}, once more
  \item \emph{Say it:} \emph{davākhānyāt sakāḷpāsūn rugṇāñcī rāṅg hotī.} — There was a queue of patients at the clinic from the morning.
  \item \emph{Recall:} say the Marathi for to roast, then the Marathi for to scratch, then say \emph{rugṇa} again
\end{itemize}

\begin{tcolorbox}[breakable,colback=teal!4,colframe=teal!35!black,title={Before you move on}]
What does \mr{रुग्ण} mean? (\textquotedblleft{}A sick person; a patient\textquotedblright{}.) Say it once more.
\end{tcolorbox}
\section[\texorpdfstring{\mr{परिचारिका} (paricārikā)}{परिचारिका (paricārikā)}]{\mr{परिचारिका} (paricārikā) — a nurse}

\label{lesson:MR-C488-paricarika}



\begin{quote}
Before the new one: say the Marathi for to scratch, then the Marathi for a sick person.
\end{quote}

\subsection*{You'll want to know: \mr{परिचारिका}}
\textbf{\mr{परिचारिका}} — \emph{paricārikā} — \textquotedblleft{}a nurse\textquotedblright{}.

\begin{quote}
\mr{रात्री} \mr{परिचारिकेने} \mr{आजोबांची} \mr{काळजी} \mr{घेतली}.
\end{quote}

\begin{quote}
\emph{rātrī paricārikene ājobāñcī kāḷjī ghetlī.}
\end{quote}

\begin{quote}
At night a nurse looked after Grandfather.
\end{quote}

Feminine.

Feminine; oblique \mr{परिचारिके}-.

\begin{cousinweb}
\textbf{\mr{परिचारिका}} (\emph{paricārikā}) is a tatsama, taken from Sanskrit \emph{paricārikā}, \textquotedblleft{}a female attendant\textquotedblright{} (Monier-Williams).
\end{cousinweb}

\begin{grammarlens}[title={What you've built}]
One more word for describing what you lived and felt.
\end{grammarlens}

\subsection*{Your turn}
\begin{itemize}
\raggedright
  \item \emph{Say it:} \emph{paricārikā}
  \item \emph{Say it:} \emph{paricārikā}, once more
  \item \emph{Say it:} \emph{rātrī paricārikene ājobāñcī kāḷjī ghetlī.} — At night a nurse looked after Grandfather.
  \item \emph{Recall:} say the Marathi for to scratch, then the Marathi for a sick person, then say \emph{paricārikā} again
\end{itemize}

\begin{tcolorbox}[breakable,colback=teal!4,colframe=teal!35!black,title={Before you move on}]
What does \mr{परिचारिका} mean? (\textquotedblleft{}A nurse\textquotedblright{}.) Say it once more.
\end{tcolorbox}
\section[\texorpdfstring{\mr{वैद्य} (vaidya)}{वैद्य (vaidya)}]{\mr{वैद्य} (vaidya) — a physician}

\label{lesson:MR-C488-vaidya}



\begin{quote}
Before the new one: say the Marathi for a sick person, then the Marathi for a nurse.
\end{quote}

\subsection*{You'll want to know: \mr{वैद्य}}
\textbf{\mr{वैद्य}} — \emph{vaidya} — \textquotedblleft{}a physician\textquotedblright{}.

\begin{quote}
\mr{गावातल्या} \mr{वैद्यांनी} \mr{आजीची} \mr{नाडी} \mr{पाहिली}.
\end{quote}

\begin{quote}
\emph{gāvātlyā vaidyānnī ājīcī nāḍī pāhilī.}
\end{quote}

\begin{quote}
The village physician felt Grandmother's pulse.
\end{quote}

Masculine.

Traditionally an Ayurvedic doctor; \mr{दंतवैद्य}, a dentist.

\begin{cousinweb}
\textbf{\mr{वैद्य}} (\emph{vaidya}) is a tatsama, taken from Sanskrit \emph{vaidya}, \textquotedblleft{}learned; a physician\textquotedblright{} (CDIAL 12140).
\end{cousinweb}

\begin{grammarlens}[title={What you've built}]
One more word for describing what you lived and felt.
\end{grammarlens}

\subsection*{Your turn}
\begin{itemize}
\raggedright
  \item \emph{Say it:} \emph{vaidya}
  \item \emph{Say it:} \emph{vaidya}, once more
  \item \emph{Say it:} \emph{gāvātlyā vaidyānnī ājīcī nāḍī pāhilī.} — The village physician felt Grandmother's pulse.
  \item \emph{Recall:} say the Marathi for a sick person, then the Marathi for a nurse, then say \emph{vaidya} again
\end{itemize}

\begin{tcolorbox}[breakable,colback=teal!4,colframe=teal!35!black,title={Before you move on}]
What does \mr{वैद्य} mean? (\textquotedblleft{}A physician\textquotedblright{}.) Say it once more.
\end{tcolorbox}
\section[\texorpdfstring{\mr{निदान} (nidān)}{निदान (nidān)}]{\mr{निदान} (nidān) — a diagnosis}

\label{lesson:MR-C488-nidan}



\begin{quote}
Before the new one: say the Marathi for a nurse, then the Marathi for a physician.
\end{quote}

\subsection*{You'll want to know: \mr{निदान}}
\textbf{\mr{निदान}} — \emph{nidān} — \textquotedblleft{}a diagnosis\textquotedblright{}.

\begin{quote}
\mr{रक्त} \mr{तपासल्यानंतर} \mr{निदान} \mr{झालं}.
\end{quote}

\begin{quote}
\emph{rakta tapāslyānantar nidān jhālaṁ.}
\end{quote}

\begin{quote}
The diagnosis was made after a blood test.
\end{quote}

Neuter.

\mr{निदान} \mr{करणे}, to diagnose. As an adverb \mr{निदान} means at least.

\begin{cousinweb}
\textbf{\mr{निदान}} (\emph{nidān}) is a tatsama, taken from Sanskrit \emph{nidāna}, \textquotedblleft{}band; primary cause\textquotedblright{} (CDIAL 7196).
\end{cousinweb}

\begin{grammarlens}[title={What you've built}]
One more word for describing what you lived and felt.
\end{grammarlens}

\subsection*{Your turn}
\begin{itemize}
\raggedright
  \item \emph{Say it:} \emph{nidān}
  \item \emph{Say it:} \emph{nidān}, once more
  \item \emph{Say it:} \emph{rakta tapāslyānantar nidān jhālaṁ.} — The diagnosis was made after a blood test.
  \item \emph{Recall:} say the Marathi for a nurse, then the Marathi for a physician, then say \emph{nidān} again
\end{itemize}

\begin{tcolorbox}[breakable,colback=teal!4,colframe=teal!35!black,title={Before you move on}]
What does \mr{निदान} mean? (\textquotedblleft{}A diagnosis\textquotedblright{}.) Say it once more.
\end{tcolorbox}
\section[\texorpdfstring{\mr{उपचार} (upcār)}{उपचार (upcār)}]{\mr{उपचार} (upcār) — treatment}

\label{lesson:MR-C488-upcar}



\begin{quote}
Before the new one: say the Marathi for a physician, then the Marathi for a diagnosis.
\end{quote}

\subsection*{You'll want to know: \mr{उपचार}}
\textbf{\mr{उपचार}} — \emph{upcār} — \textquotedblleft{}treatment\textquotedblright{}.

\begin{quote}
\mr{योग्य} \mr{उपचारांनंतर} \mr{तो} \mr{बरा} \mr{झाला}.
\end{quote}

\begin{quote}
\emph{yogya upcārānnantar to barā jhālā.}
\end{quote}

\begin{quote}
After proper treatment he got well.
\end{quote}

Masculine.

Often plural: \mr{उपचार} \mr{करणे}, to treat; \mr{उपचार} \mr{घेणे}, to have treatment.

\begin{cousinweb}
\textbf{\mr{उपचार}} (\emph{upcār}) is a tatsama, taken from Sanskrit \emph{upacāra}, \textquotedblleft{}approach, service\textquotedblright{} (CDIAL 2163).
\end{cousinweb}

\begin{grammarlens}[title={What you've built}]
One more word for describing what you lived and felt.
\end{grammarlens}

\subsection*{Your turn}
\begin{itemize}
\raggedright
  \item \emph{Say it:} \emph{upcār}
  \item \emph{Say it:} \emph{upcār}, once more
  \item \emph{Say it:} \emph{yogya upcārānnantar to barā jhālā.} — After proper treatment he got well.
  \item \emph{Recall:} say the Marathi for a physician, then the Marathi for a diagnosis, then say \emph{upcār} again
\end{itemize}

\begin{tcolorbox}[breakable,colback=teal!4,colframe=teal!35!black,title={Before you move on}]
What does \mr{उपचार} mean? (\textquotedblleft{}Treatment\textquotedblright{}.) Say it once more.
\end{tcolorbox}
